\section{Worked representation and selector calculations}\label{sec:worked_calculations}

The preceding arguments use a small number of mechanisms repeatedly. This section expands those mechanisms in examples that can be checked without automorphic input. The examples are finite representation models. An incompatible local model is not an arithmetic counterexample to Strong Artin.

\subsection{Why the prime cyclic Clifford alternative is exact}

\begin{proposition}[The two prime cyclic cases]\label{prop:worked_clifford}
Let \(G\) be finite, let \(H\triangleleft G\) with \(G/H\cong C_p\), and let \(V\) be irreducible over \(\CC\). Exactly one of the following occurs:
\begin{enumerate}
\item \(V|_H\) is irreducible and invariant under \(G\), and \(V\) is one of its \(p\) quotient character extensions.
\item \(V|_H\) is a sum of \(p\) pairwise nonisomorphic conjugate irreducible representations, and \(V\) is induced from any one of them.
\end{enumerate}
\end{proposition}

\begin{proof}
Complete reducibility for the finite subgroup gives an isotypic decomposition of \(V|_H\). The group \(G\) permutes these isotypic spaces by conjugating their types. More than one orbit would give a proper nonzero \(G\) stable direct sum, contrary to irreducibility. Thus there is one orbit. Since conjugation by \(H\) fixes isomorphism classes of \(H\) representations, the permutation action factors through \(C_p\). Its transitive orbit has size one or \(p\).

In the size \(p\) case, choose one isotypic space \(U\). Its translates are pairwise distinct isotypic spaces, with direct sum \(V\). The stabilizer of its type is exactly \(H\). If \(U\) had a proper nonzero \(H\) submodule, the direct sum of its \(p\) translates would be a proper nonzero \(G\) submodule of \(V\). Hence \(U\) is irreducible. The natural map from \(\Ind_H^G U\) to \(V\), sending each coset copy to its translated space, is an isomorphism. Its restriction is precisely the stated orbit sum.

In the size one case, write \(V|_H\cong\tau\otimes M\), where \(H\) acts trivially on the multiplicity space \(M\). The stable type \(\tau\) extends by \cref{lem:stable_extension}. Remove that extension action from the action on \(V\). The remaining action on \(M\) factors through \(C_p\). If \(M\) had a proper quotient stable subspace, its tensor product with \(\tau\) would give a proper \(G\) submodule of \(V\). Thus \(M\) is irreducible as a complex representation of \(C_p\). Every such irreducible is one dimensional. Therefore \(V|_H\cong\tau\), with multiplicity one, and the remaining quotient character is its extension label. The torsor theorem gives the \(p\) possible extensions. This proves both alternatives and excludes an intermediate multiplicity case.
\end{proof}

For a solvable finite group, a subnormal composition series has prime cyclic factors. Applied recursively to the constituents encountered along the series, this proposition separates induction steps from invariant extension steps. It does not say that every irreducible representation is induced all the way from a linear character. An invariant restriction can survive at positive rank. That surviving case is the attachment gate used in the main text.

For \(S_3\) and \(H=A_3\), the standard two dimensional representation lies in the moving case. A three cycle has the two nontrivial cube roots as eigenvalues. Restriction is the sum of the two nontrivial \(C_3\) characters. Induction from either one gives the standard representation. Its sign twist is isomorphic to itself. This self twist is allowed precisely because the restricted representation is reducible. It does not contradict freeness of the torsor whose restriction is irreducible.

\subsection{A rank three coefficient test that determinant misses}

Let \(\omega\) be a primitive cube root of unity and put
\[
 A=\operatorname{diag}(1,1,\omega).
\]
Expanding directly gives
\[
 \det(TI-A)
 =T^3-(2+\omega)T^2+(1+2\omega)T-\omega.
\]
The three elementary symmetric coefficients are
\[
 e_1=2+\omega,\qquad e_2=1+2\omega,\qquad e_3=\omega.
\]
For scalar twist \(\omega^jA\), they become
\[
 \omega^j(2+\omega),\qquad
 \omega^{2j}(1+2\omega),\qquad
 \omega^{3j}\omega=\omega.
\]
The determinant is constant for \(j=0,1,2\), but the first coefficient distinguishes the three labels because \(2+\omega\ne0\). Thus rank divisibility causes blindness of the determinant, rather than blindness of every coefficient.

Now compare the balanced matrix
\[
 A_{\mathrm{bal}}=\operatorname{diag}(1,\omega,\omega^2).
\]
Its polynomial is \(T^3-1\). Multiplying all eigenvalues by \(\omega\) merely permutes their multiset. Every lower coefficient vanishes, and the determinant is fixed. Its full characteristic polynomial is therefore blind to the label. This is an actual local rotational stabilizer. A claim that one characteristic polynomial always detects a quotient twist would fail on this example.

At quotient prime three and rank two, the determinant changes by \(\omega^{2j}\). If a required determinant ratio is \(\omega^b\), the unique candidate satisfies \(2j=b\) modulo three, so \(j=2b\). This solves a finite selection problem. It still does not show that every other coefficient of the candidate matches the target.

The all local comparison includes infinity. For a real place with complex quadratic extension, the quotient is \(C_2\). Let the target finite parameter have sign multiplicities \((3,0)\), and let a second parameter have multiplicities \((1,2)\). Both have determinant sign positive. Twisting the second by sign gives \((2,1)\), which still does not equal \((3,0)\). Its exact attachment set is empty. This is a finite local model showing why determinant selection, even with quotient prime coprime to rank, does not discharge the remaining infinity condition.

\subsection{Ordered multiplicities and a common shift}

Take \(G=C_2\times C_3\), let \(H=C_2\), and let \(\xi\) be the quotient character taking a generator of \(C_3\) to \(\omega\). Let \(\tau\) be the nontrivial character of \(H\), extended with trivial quotient action. Consider a semisimple representation with spectrum
\[
 P_{\tau,V}(z)=2+z.
\]
Its quotient multiplicity vector is \((2,1,0)\). Twisting by \(\xi\) gives
\[
 zP_{\tau,V}(z)=2z+z^2,
\]
whose vector is \((0,2,1)\). Twisting by \(\xi^2\) gives \((1,0,2)\). These are ordered vectors. Knowing that each has total multiplicity three does not retain the extension label.

Suppose a second stable \(H\) type \(\sigma\) occurs with spectrum one. If the target has spectra \(2z+z^2\) for \(\tau\) and \(z^2\) for \(\sigma\), the first block demands shift one and the second demands shift two. The two restrictions to \(H\) have the same multiplicities. Each block separately is a quotient twist of the matching source block. The full representation is not one quotient twist, because the required shifts disagree. This is the reason the spectrum criterion quantifies over one common \(j\), rather than an independent \(j_\tau\) for each type.

For a balanced multiplicity spectrum \(1+z+z^2\), both nonzero Fourier modes vanish:
\[
 1+\omega+\omega^2=0,\qquad
 1+\omega^2+\omega^4=0.
\]
The spectrum \(3\) has the same total multiplicity, but both nonzero modes equal three. Comparing only ratios for which the source mode is nonzero would omit both decisive comparisons. The zero support clause in \cref{def:local_spectral_obstruction} is therefore necessary. A zero source mode requires a zero target mode. It is not an instruction to skip that mode.

Changing the reference extension of \(\tau\) relabels both spectra by the same rotation. The attachment set does not change. Likewise changing the chosen generator of the quotient rescales all labels by an invertible element of \(\ZZ/p\ZZ\). These coordinate choices affect the displayed polynomials, but not the existence of an isomorphism after one common quotient twist.

For prime cyclic quotients a nonempty exact alignment set is either a singleton or all labels. Hence two nonempty exact singleton conditions \(\{1\}\) and \(\{2\}\) are incompatible, while \(\{1\}\) together with all labels is compatible. Coarse coefficient filters may produce other finite subsets, but those are not automatically exact extension spectrum sets. Keeping that distinction prevents a visual or computational filter from being assigned a stronger meaning than its construction supports.

\subsection{A trace encoding, with both hypothesis failures exposed}

Let
\[
 \Gamma=\{\operatorname{diag}(\omega^a,\omega^{-a}):a=0,1,2\}.
\]
The trace alphabet is \(\{2,-1\}\). The image algebra has dimension two, with basis the identity and \(\gamma_1=\operatorname{diag}(\omega,\omega^2)\). The trace coordinate map from \cref{thm:finite_trace_image_bound} gives
\[
 \Phi(\gamma_0)=(2,-1),\qquad
 \Phi(\gamma_1)=(-1,-1),\qquad
 \Phi(\gamma_2)=(-1,2).
\]
The tuples are distinct. The sharper algebra bound is \(3\le2^2=4\), while the rank only bound is \(3\le2^{2^2}=16\). The latter is deliberately rough. Its role is a uniform bound independent of the coefficient prime, not an efficient classification of the group.

Semisimplicity cannot be omitted. Over \(\CC\), the matrices
\[
 \begin{pmatrix}1&a\\0&1\end{pmatrix},\qquad a\in\ZZ,
\]
form an infinite group with the single trace value two. The invariant first coordinate line has no invariant complementary line. The trace pairing does not see the upper right nilpotent direction. This is exactly the step that fails without semisimplicity.

The characteristic restriction also has a visible purpose. In characteristic two, choose a primitive cube root \(\omega\) and the scalar group
\[
 \{I_2,\omega I_2,\omega^2 I_2\}.
\]
It is a faithful semisimple representation, as a direct sum of two identical one dimensional characters. Every trace is zero because its multiplicity two vanishes in the coefficient field. The alphabet has size one, but the group has size three. The would be bound \(3\le1^{4}\) is false. In the proof, the block pairing multiplier \(m_i=2\) is zero. Requiring \(\ell>n\) prevents precisely this multiplicity cancellation.

In characteristic zero the finite trace theorem itself is enough to prove finiteness once global semisimplicity and a finite image trace alphabet are known. The l adic part of the work is therefore not another mysterious finite image principle. It is the arithmetic passage from density one Frobenius data to the trace set of the full compact image, followed by this explicit algebraic encoding.

\subsection{A fixed alphabet and what a separating prime actually does}

For an elementary example, compare the integral polynomials
\[
 P(T)=T^2-2T+1,\qquad Q(T)=T^2+2T+1.
\]
Their middle coefficients differ by four. Reduction modulo two makes them equal. Reduction modulo any odd rational prime keeps them distinct. A separating prime is a coefficient prime with this latter injectivity property. It does not say which polynomial is realized by a global Galois representation. It says that, once a realization gives one of the alphabet polynomials modulo that prime, a different alphabet entry cannot imitate it.

For a cyclotomic alphabet over a number field, a nonzero algebraic integer difference has only finitely many prime ideal divisors. Choosing one nonzero coefficient difference for every pair therefore produces the finite bad set. Distinct roots can reduce to the same root at some bad prime, and repeated root multiplicities can become inseparable. The argument avoids those accidents by testing the polynomial alphabet, rather than presuming that each individual eigenvalue is already distinguished.

The large residual theorem uses two alphabets. One is the supplied automorphic alphabet. The other contains every degree \(n\) polynomial whose roots have orders at most \(B=t^{n^2}\). A field containing \(\mu_{\operatorname{lcm}(1,\ldots,B)}\) contains all their coefficients. Both the field and the exceptional coefficient primes are fixed before choosing a residual realization. This order matters. A proof that chose a separation field from the coefficients of an as yet unknown lift and then called the excluded primes a predetermined set would leave a circular quantifier.

The resulting bound is not claimed to be practical. Even \(t\) and \(n\) of moderate size make \(B\) and the universal cyclotomic field enormous. The theorem establishes an existence implication with a finite exceptional set. A useful computational implementation for a particular finite image may exploit its actual exponent and much smaller image algebra, but that is an additional optimization and must retain the same separation logic.

\section{Statement map and the remaining arithmetic obligation}\label{sec:statement_map}

The reader edition has three layers of evidence. The finite algebra proofs in the main text and local appendix are written out. Classical automorphic, local Langlands, and representation realization results are imported with their source hypotheses. The formal archive records exactly encoded algebraic theorems and historical build evidence. These layers support one another without being interchangeable.

\begin{longtable}{>{\raggedright\arraybackslash}p{.24\linewidth}>{\raggedright\arraybackslash}p{.33\linewidth}>{\raggedright\arraybackslash}p{.34\linewidth}}
\toprule
Statement group & What the proof supplies & What remains an input\\
\midrule
\endfirsthead
\toprule
Statement group & What the proof supplies & What remains an input\\
\midrule
\endhead
Cyclic extension torsor & All extensions of one invariant irreducible type differ by a unique quotient character & An automorphic attachment to select the matching extension\\
Coefficient and multiplicity selectors & Exact rotation, stabilizer, and common shift tests & Actual local data and compatibility of the selected global label\\
Finite Satake descent & A bounded root of unity alphabet outside a finite exceptional set & The known lower automorphic correspondence and base change compatibility\\
Trace encoding & Finite image and the bound under semisimplicity and characteristic hypotheses & A representation whose trace alphabet has the required form\\
Separating residual recovery & Exact weak characteristic zero attachment from a supplied good residual realization & Existence of that residual realization and remaining local checks\\
Large residual recovery & A finite ordinary lift and exact alphabet recovery at sufficiently large good primes & The residual representation and classical cross characteristic lifting theorem\\
Quotient prime congruence & Equal reduced unramified polynomial systems with collapsed labels & A suitable weight zero lifting theorem with target local control\\
Semisimple l adic rigidity & Finite image from continuous finitely ramified realization and finite Frobenius traces & Existence and exact polynomial normalization of that realization\\
Cyclic cohomology and syzygies & Algebraic spaces, redundancy relations, and conditional local obstruction criteria & Arithmetic Selmer comparison and the required terminal obstruction vanishing\\
Certified matrix Hensel sector & Unique actual local root from one finite test when \(m>s\) and the two sided inverse certificate holds & A certificate, singular sector analysis, and arithmetic or global realization\\
Solvable recursion & A complete implication from gates at every invariant step & Uniformly supplying those gates over the encountered fields and ranks\\
\bottomrule
\end{longtable}

The historical paper contains additional theorem statements, corollaries, and small example verifiers. Its source and audit are retained as evidence, with the current statement map recording the narrowed claims used in this edition. The older global semisimplicity, Fourier zero support, coefficient field, and exceptional place issues have been corrected in the reader text. The later formal checkpoints are not retroactive certificates of every classical theorem or every prose argument in that historical paper.

The remaining global question can now be stated without a vague final step. At an invariant cyclic restriction, the automorphic descent fibre is known to exist. Can one prove that a descent has a continuous Galois realization with the exact normalized polynomial data, or lift the canonical quotient prime congruence to the prescribed target with its local conditions? If the answer supplies only almost everywhere polynomials, the exceptional finite and archimedean conditions must still be addressed for the all local version.

This is an existence and compatibility obligation. A finite trace alphabet controls a representation after it exists. A syzygy certificate removes redundant equations after the deformation comparison is supplied. A Hensel certificate solves a specified local matrix equation in its certified ball. None of those conclusions creates the global automorphic to Galois arrow by itself. Their value is that the resulting arrow, when constructed, will arrive at a tightly constrained and independently checkable destination.
